\subsection{Biểu đồ tuần tự}

Phần này trình bày các biểu đồ tuần tự (\textit{Sequence Diagram}) mô tả luồng tương tác chi tiết giữa các thành phần trong hệ thống cho từng chức năng đã được đặc tả ở mục \ref{tab:uc-xac-thuc-nguoi-dung} và các mục đặc tả tiếp theo. Mỗi biểu đồ thể hiện trình tự trao đổi thông điệp qua các tầng: tác nhân, BFF Gateway, các microservice nghiệp vụ và các thành phần lưu trữ (MongoDB, Redis, Blockchain).

Bước \textit{Xác thực và Kiểm tra quyền} xuất hiện lặp lại trong hầu hết các luồng yêu cầu xác thực; do đó được tách thành một biểu đồ tham chiếu dùng chung (Hình~\ref{fig:seq-ref-auth}) và được tham chiếu dưới dạng \texttt{ref} trong các biểu đồ còn lại nhằm tránh trùng lặp. Biểu đồ này mô tả cách BFF Gateway kiểm tra access token, trích xuất định danh người dùng và đối chiếu vai trò yêu cầu trước khi cho phép xử lý tiếp. Các nhánh lỗi như thiếu token, token hết hạn hoặc vai trò không phù hợp đều bị chặn ngay tại tầng BFF.

\begin{figure}[H]
    \centering
    \includegraphics[width=0.5\textwidth]{Diagrams/Sequence/ref_auth_check.pdf}
    \caption{Biểu đồ tuần tự tham chiếu ``Xác thực và Kiểm tra quyền''}
    \label{fig:seq-ref-auth}
\end{figure}

\subsubsection{Xác thực người dùng}

Biểu đồ tuần tự của chức năng xác thực người dùng mô tả ba thao tác trong vòng đời phiên làm việc: đăng nhập, làm mới phiên và đăng xuất. BFF Gateway tiếp nhận yêu cầu từ người dùng, sau đó chuyển tiếp tới Identity Service để kiểm tra tài khoản, mật khẩu và trạng thái kích hoạt. Redis được sử dụng để lưu, xoay vòng và thu hồi refresh token nhằm bảo đảm phiên đăng nhập có thể bị chấm dứt một cách chủ động.

\begin{figure}[H]
    \centering
    \includegraphics[width=0.9\textwidth, keepaspectratio]{Diagrams/Sequence/XacThucNguoiDung.pdf}
    \caption{Biểu đồ tuần tự ``Xác thực người dùng''}
    \label{fig:seq-xac-thuc-nguoi-dung}
\end{figure}

\subsubsection{Tạo tài khoản người dùng}

Biểu đồ tuần tự của chức năng tạo tài khoản người dùng mô tả cách quản trị viên gửi yêu cầu tạo một hoặc nhiều tài khoản qua BFF Gateway. Identity Service kiểm tra dữ liệu đầu vào, phát hiện email trùng lặp, băm mật khẩu và ghi bản ghi người dùng vào MongoDB. Với trường hợp tạo hàng loạt, luồng trả về tổng hợp kết quả để quản trị viên biết các tài khoản đã được tạo hoặc các bản ghi cần chỉnh sửa.

\begin{figure}[H]
    \centering
    \includegraphics[width=0.9\textwidth, keepaspectratio]{Diagrams/Sequence/TaoDanhSachNguoiDung.pdf}
    \caption{Biểu đồ tuần tự ``Tạo tài khoản người dùng''}
    \label{fig:seq-tao-danh-sach-nguoi-dung}
\end{figure}

\subsubsection{Tìm kiếm người dùng}

Biểu đồ tuần tự này mô tả quá trình quản trị viên tìm kiếm tài khoản theo từ khóa, vai trò, trạng thái và thông tin phân trang. Sau khi BFF Gateway xác thực quyền quản trị, Identity Service xây dựng truy vấn lọc trên cơ sở dữ liệu người dùng. Kết quả trả về gồm danh sách bản ghi theo trang và tổng số bản ghi phù hợp với điều kiện lọc.

\begin{figure}[H]
    \centering
    \includegraphics[width=0.9\textwidth, keepaspectratio]{Diagrams/Sequence/TimKiemNguoiDung.pdf}
    \caption{Biểu đồ tuần tự ``Tìm kiếm người dùng''}
    \label{fig:seq-tim-kiem-nguoi-dung}
\end{figure}

\subsubsection{Quản lý trạng thái tài khoản}

Biểu đồ tuần tự này mô tả thao tác kích hoạt hoặc vô hiệu hóa tài khoản người dùng. Sau khi yêu cầu được xác thực quyền quản trị tại BFF Gateway, Identity Service truy vấn tài khoản mục tiêu và cập nhật trạng thái trong cơ sở dữ liệu. Trạng thái mới được trả về cho quản trị viên, đồng thời là cơ sở để hệ thống từ chối đăng nhập hoặc bỏ phiếu đối với tài khoản đã bị vô hiệu hóa.

\begin{figure}[H]
    \centering
    \includegraphics[width=0.9\textwidth, keepaspectratio]{Diagrams/Sequence/QuanLyTrangThaiTaiKhoan.pdf}
    \caption{Biểu đồ tuần tự ``Quản lý trạng thái tài khoản''}
    \label{fig:seq-quan-ly-trang-thai-tai-khoan}
\end{figure}

\subsubsection{Tạo cuộc bầu cử}

Biểu đồ tuần tự của chức năng tạo cuộc bầu cử mô tả luồng quản trị viên khởi tạo một cuộc bầu cử mới ở trạng thái PENDING. Coordinator Service tiếp nhận dữ liệu, yêu cầu Identity Service kiểm tra danh sách ứng viên ban đầu và các ràng buộc cấu hình như số ứng viên tối đa được chọn. Khi dữ liệu hợp lệ, cuộc bầu cử được ghi vào MongoDB của Coordinator Service và trả về cho quản trị viên.

\begin{figure}[H]
    \centering
    \includegraphics[width=0.9\textwidth, keepaspectratio]{Diagrams/Sequence/TaoCuocBauCu.pdf}
    \caption{Biểu đồ tuần tự ``Tạo cuộc bầu cử''}
    \label{fig:seq-tao-cuoc-bau-cu}
\end{figure}

\subsubsection{Tìm kiếm cuộc bầu cử}

Biểu đồ tuần tự này thể hiện cách quản trị viên tra cứu danh sách cuộc bầu cử theo từ khóa, trạng thái, khoảng thời gian và phân trang. BFF Gateway kiểm tra quyền truy cập trước khi chuyển yêu cầu tới Coordinator Service. Coordinator Service xây dựng truy vấn trên cơ sở dữ liệu cuộc bầu cử và trả về danh sách phù hợp để phục vụ các thao tác quản trị tiếp theo.

\begin{figure}[H]
    \centering
    \includegraphics[width=0.9\textwidth, keepaspectratio]{Diagrams/Sequence/TimKiemCuocBauCu.pdf}
    \caption{Biểu đồ tuần tự ``Tìm kiếm cuộc bầu cử''}
    \label{fig:seq-tim-kiem-cuoc-bau-cu}
\end{figure}

\subsubsection{Xem chi tiết cuộc bầu cử}

Biểu đồ tuần tự của chức năng xem chi tiết cuộc bầu cử mô tả luồng tổng hợp thông tin cấu hình, danh sách ứng viên và danh sách cử tri. Coordinator Service đọc bản ghi cuộc bầu cử trong MongoDB, sau đó phối hợp với Identity Service để lấy hồ sơ chi tiết của các tài khoản liên quan. Nếu cuộc bầu cử không tồn tại, luồng dừng tại Coordinator Service và trả về lỗi cho quản trị viên.

\begin{figure}[H]
    \centering
    \includegraphics[width=0.9\textwidth, keepaspectratio]{Diagrams/Sequence/XemChiTietCuocBauCu.pdf}
    \caption{Biểu đồ tuần tự ``Xem chi tiết cuộc bầu cử''}
    \label{fig:seq-xem-chi-tiet-cuoc-bau-cu}
\end{figure}

\subsubsection{Quản lý ứng viên trong cuộc bầu cử}

Biểu đồ tuần tự này mô tả thao tác thêm hoặc xóa ứng viên khỏi một cuộc bầu cử cụ thể. Coordinator Service trước hết kiểm tra cuộc bầu cử có tồn tại và còn ở trạng thái PENDING hay không, nhằm tránh thay đổi danh sách ứng viên sau khi cuộc bầu cử đã bắt đầu. Identity Service được gọi để xác minh các tài khoản ứng viên trước khi Coordinator Service cập nhật danh sách trong cơ sở dữ liệu.

\begin{figure}[H]
    \centering
    \includegraphics[width=0.9\textwidth, keepaspectratio]{Diagrams/Sequence/QuanLyUngVienTrongBauCu.pdf}
    \caption{Biểu đồ tuần tự ``Quản lý ứng viên trong cuộc bầu cử''}
    \label{fig:seq-quan-ly-ung-vien-trong-bau-cu}
\end{figure}

\subsubsection{Quản lý cử tri trong cuộc bầu cử}

Biểu đồ tuần tự của chức năng quản lý cử tri mô tả cách quản trị viên thêm hoặc xóa cử tri khỏi danh sách được phép tham gia một cuộc bầu cử. Coordinator Service kiểm tra trạng thái PENDING của cuộc bầu cử trước khi cho phép thay đổi, vì danh sách cử tri quyết định quyền bỏ phiếu. Identity Service xác minh danh sách tài khoản cử tri, sau đó Coordinator Service cập nhật quan hệ cử tri trong MongoDB.

\begin{figure}[H]
    \centering
    \includegraphics[width=0.9\textwidth, keepaspectratio]{Diagrams/Sequence/QuanLyCuTriTrongBauCu.pdf}
    \caption{Biểu đồ tuần tự ``Quản lý cử tri trong cuộc bầu cử''}
    \label{fig:seq-quan-ly-cu-tri-trong-bau-cu}
\end{figure}

\subsubsection{Điều hành cuộc bầu cử}

Biểu đồ tuần tự này mô tả hai mốc chính trong vòng đời cuộc bầu cử: bắt đầu và đóng cuộc bầu cử. Khi bắt đầu, Coordinator Service kiểm tra điều kiện cấu hình và yêu cầu các Signing Node sinh khóa công khai tập thể EC-Schnorr trước khi chuyển trạng thái sang ACTIVE. Khi đóng, hệ thống tổng hợp các cam kết phiếu, sinh gốc Merkle và ghi cam kết này lên blockchain qua Chainlaunch Gateway.

\begin{figure}[H]
    \centering
    \includegraphics[width=0.9\textwidth, keepaspectratio]{Diagrams/Sequence/DieuHanhCuocBauCu.pdf}
    \caption{Biểu đồ tuần tự ``Điều hành cuộc bầu cử''}
    \label{fig:seq-dieu-hanh-cuoc-bau-cu}
\end{figure}

\subsubsection{Xem danh sách phiếu trong cuộc bầu cử}

Biểu đồ tuần tự của chức năng xem danh sách phiếu mô tả luồng quản trị viên giám sát các phiếu đã nộp trong một cuộc bầu cử. Coordinator Service truy vấn các bản ghi phiếu theo định danh cuộc bầu cử, kèm thông tin phân trang và các trường công khai như cam kết phiếu mù, thời điểm tạo, tham chiếu giao dịch blockchain. Luồng này không trả về lựa chọn ứng viên hay bất kỳ dữ liệu bí mật nào của cử tri.

\begin{figure}[H]
    \centering
    \includegraphics[width=0.9\textwidth, keepaspectratio]{Diagrams/Sequence/XemDanhSachPhieu.pdf}
    \caption{Biểu đồ tuần tự ``Xem danh sách phiếu trong cuộc bầu cử''}
    \label{fig:seq-xem-danh-sach-phieu}
\end{figure}

\subsubsection{Xem danh sách cuộc bầu cử của cử tri}

Biểu đồ tuần tự này thể hiện luồng cử tri xem các cuộc bầu cử mà mình thuộc danh sách tham gia. Sau khi BFF Gateway xác thực token và vai trò cử tri, Coordinator Service tra cứu quan hệ giữa cử tri và cuộc bầu cử trong MongoDB. Kết quả trả về gồm thông tin tóm tắt từng cuộc bầu cử và trạng thái bỏ phiếu của chính cử tri.

\begin{figure}[H]
    \centering
    \includegraphics[width=0.9\textwidth, keepaspectratio]{Diagrams/Sequence/XemBauCuCuaToi.pdf}
    \caption{Biểu đồ tuần tự ``Xem danh sách cuộc bầu cử của cử tri''}
    \label{fig:seq-xem-bau-cu-cua-toi}
\end{figure}

\subsubsection{Xem chi tiết cuộc bầu cử đang tham gia}

Biểu đồ tuần tự của chức năng này mô tả cách cử tri xem thông tin chi tiết của một cuộc bầu cử trước khi bỏ phiếu. Coordinator Service kiểm tra cử tri có thuộc danh sách hợp lệ của cuộc bầu cử hay không, sau đó lấy cấu hình cuộc bầu cử và danh sách ứng viên. Identity Service được sử dụng để bổ sung hồ sơ ứng viên, còn trạng thái bỏ phiếu của cử tri được trả về cùng dữ liệu chi tiết.

\begin{figure}[H]
    \centering
    \includegraphics[width=0.9\textwidth, keepaspectratio]{Diagrams/Sequence/XemChiTietBauCuThamGia.pdf}
    \caption{Biểu đồ tuần tự ``Xem chi tiết cuộc bầu cử đang tham gia''}
    \label{fig:seq-xem-chi-tiet-bau-cu-tham-gia}
\end{figure}

\subsubsection{Bỏ phiếu}

Biểu đồ tuần tự của chức năng bỏ phiếu mô tả ba bước chính: khởi tạo phiên bỏ phiếu, yêu cầu ký phiếu mù tập thể và nộp cam kết phiếu mù. Coordinator Service phối hợp với Redis để duy trì phiên bỏ phiếu tạm thời và với các Signing Node để sinh chữ ký EC-Schnorr từng phần. Sau khi cử tri nộp cam kết phiếu mù, hệ thống ghi nhận phiếu trong cơ sở dữ liệu và blockchain mà không tiết lộ lựa chọn ứng viên.

\begin{figure}[H]
    \centering
    \includegraphics[height=0.9\textheight, keepaspectratio]{Diagrams/Sequence/BoPhieu.pdf}
    \caption{Biểu đồ tuần tự ``Bỏ phiếu''}
    \label{fig:seq-bo-phieu}
\end{figure}

\subsubsection{Tiết lộ phiếu sau khi cuộc bầu cử đóng}

Biểu đồ tuần tự này mô tả bước công khai phiếu sau khi cuộc bầu cử đã đóng. Reveal-Vote Service nhận yêu cầu từ khách ẩn danh, kiểm tra trạng thái cuộc bầu cử, chuẩn hóa danh sách ứng viên và xác minh chữ ký EC-Schnorr bằng khóa công khai tập thể. Khi hợp lệ, phiếu tiết lộ được ghi lên blockchain và cơ sở dữ liệu; nếu tất cả phiếu đã được tiết lộ, hệ thống chuyển cuộc bầu cử sang trạng thái hoàn tất.

\begin{figure}[H]
    \centering
    \includegraphics[width=0.9\textwidth, keepaspectratio]{Diagrams/Sequence/TietLoPhieu.pdf}
    \caption{Biểu đồ tuần tự ``Tiết lộ phiếu sau khi cuộc bầu cử đóng''}
    \label{fig:seq-tiet-lo-phieu}
\end{figure}

\subsubsection{Xác minh phiếu bầu}

Biểu đồ tuần tự của chức năng xác minh phiếu bầu mô tả quy trình khách dùng biên lai để đối chiếu một phiếu cụ thể. Coordinator Service kiểm tra bản ghi trong cơ sở dữ liệu, đối chiếu dữ liệu trên blockchain và, khi cuộc bầu cử đã đóng, kiểm tra thêm bằng chứng Merkle. Kết quả cuối cùng chỉ được coi là hợp lệ khi các tầng kiểm tra đều xác nhận khớp.

\begin{figure}[H]
    \centering
    \includegraphics[width=0.9\textwidth, keepaspectratio]{Diagrams/Sequence/XacMinhPhieuBau.pdf}
    \caption{Biểu đồ tuần tự ``Xác minh phiếu bầu''}
    \label{fig:seq-xac-minh-phieu-bau}
\end{figure}

\subsubsection{Xem kết quả bầu cử}

Biểu đồ tuần tự này mô tả cách khách công khai xem kết quả sau khi cuộc bầu cử đã đóng hoặc hoàn tất. Reveal-Vote Service kiểm tra trạng thái cuộc bầu cử, đếm số phiếu đã tiết lộ trong cơ sở dữ liệu, truy vấn kết quả kiểm phiếu trên blockchain và lấy hồ sơ ứng viên từ Identity Service. Kết quả trả về kèm cả số liệu từ cơ sở dữ liệu và blockchain để hỗ trợ đối chiếu chéo.

\begin{figure}[H]
    \centering
    \includegraphics[width=0.9\textwidth, keepaspectratio]{Diagrams/Sequence/XemKetQuaBauCu.pdf}
    \caption{Biểu đồ tuần tự ``Xem kết quả bầu cử''}
    \label{fig:seq-xem-ket-qua-bau-cu}
\end{figure}

\subsubsection{Kiểm toán số phiếu}

Biểu đồ tuần tự của chức năng kiểm toán số phiếu mô tả luồng khách lấy các chỉ số đối chiếu giữa hệ thống nội bộ và blockchain. Reveal-Vote Service tổng hợp số phiếu đã tiết lộ từ cơ sở dữ liệu, số phiếu đã nộp từ Coordinator Service và các chỉ số kiểm toán từ chaincode. Kết quả giúp người xem kiểm tra tính nhất quán của tổng số phiếu, số phiếu đã tiết lộ và trạng thái cam kết gốc Merkle.

\begin{figure}[H]
    \centering
    \includegraphics[width=0.9\textwidth, keepaspectratio]{Diagrams/Sequence/KiemToanSoPhieu.pdf}
    \caption{Biểu đồ tuần tự ``Kiểm toán số phiếu''}
    \label{fig:seq-kiem-toan-so-phieu}
\end{figure}

\subsubsection{Sao lưu khóa bí mật phiếu}

Biểu đồ tuần tự này mô tả luồng cử tri sao lưu khóa bí mật phiếu từ thiết bị di động lên máy chủ. Ứng dụng dẫn xuất khóa mã hóa từ mã PIN, mã hóa dữ liệu cục bộ bằng AES-256-GCM rồi gửi gói envelope qua BFF Gateway tới Identity Service. Máy chủ chỉ lưu dữ liệu đã mã hóa và áp dụng thêm lớp bảo vệ khi lưu trữ, nên không nắm được mã PIN hay nội dung khóa.

\begin{figure}[H]
    \centering
    \includegraphics[width=0.9\textwidth, keepaspectratio]{Diagrams/Sequence/SaoLuuKhoa.pdf}
    \caption{Biểu đồ tuần tự ``Sao lưu khóa bí mật phiếu''}
    \label{fig:seq-sao-luu-khoa}
\end{figure}

\subsubsection{Khôi phục khóa bí mật phiếu}

Biểu đồ tuần tự của chức năng khôi phục khóa bí mật phiếu mô tả cách cử tri lấy lại bản sao lưu khi đổi thiết bị hoặc cài lại ứng dụng. Identity Service trả về gói envelope đã lưu sau khi yêu cầu được xác thực, còn việc dẫn xuất khóa và giải mã diễn ra hoàn toàn trên thiết bị. Nếu chưa có bản sao lưu hoặc mã PIN không đúng, luồng trả về lỗi tương ứng mà không làm lộ nội dung khóa.

\begin{figure}[H]
    \centering
    \includegraphics[width=0.9\textwidth, keepaspectratio]{Diagrams/Sequence/KhoiPhucKhoa.pdf}
    \caption{Biểu đồ tuần tự ``Khôi phục khóa bí mật phiếu''}
    \label{fig:seq-khoi-phuc-khoa}
\end{figure}
